\chapter{\astar Search}
\label{ch:ch04}
% Function names for pseudocode are prefixed with "Astar" to stay unique
% across the book.
\SetKwFunction{AstarSearch}{AStar}
\SetKwFunction{AstarSpaceTime}{SpaceTimeAStar}
\SetKwFunction{AstarPath}{ReconstructPath}

\begin{objectives}
\begin{itemize}
  \item Explain best-first search with $\fcost = \gcost + \hcost$ and trace \astar by hand on a small grid while keeping track of \Open and \Closed.
  \item State the exact definitions of admissible and consistent heuristics, prove that consistency implies admissibility, and prove that \astar returns an optimal path.
  \item Choose a heuristic for a 4- or 8-connected grid (Manhattan, octile, Euclidean, Chebyshev), break ties well, and predict what a weight $w$ does to cost and effort.
  \item Add time to the state: run \astar in space-time with wait actions, vertex and edge constraints, a reservation table and a correct goal test.
  \item Implement grid \astar and space-time \astar in \python, visualise the open and closed sets, and recognise the failure modes that matter in a drone planner.
\end{itemize}
\end{objectives}

% ---------------------------------------------------------------------
\section{Why this matters for a drone swarm}
\label{sec:ch04-motivation}

Picture a delivery drone that has to cross a city district represented by a
$100 \times 100$ grid of cells, a quarter of them blocked by buildings.
Dijkstra's algorithm from \cref{ch:ch03} finds the shortest route, but it does
so by settling almost every free cell: in the experiment of
\cref{sec:ch04-experiment} it expands about $7\,450$ of the $7\,500$ free cells
on such maps. The drone knows something Dijkstra ignores: the goal lies to the
north-east, and a cell far to the south-west is unlikely to be on the way. \astar
puts this knowledge into the search. With the Manhattan distance as a guide it
solves the same instances with about $750$ expansions, ten times fewer, and it
still guarantees the shortest route.

For a single query this is a convenience. For a swarm it is a necessity. The
global planner of the hybrid architecture in \cref{ch:ch24} is
conflict-based search (\cbs) or its bounded-suboptimal variant, enhanced \cbs
(\ecbs; \cref{ch:ch09,ch:ch10}), and every node of their constraint tree calls a
single-agent search for one drone under a set of constraints. A run on a few
dozen drones can easily make thousands of such calls, and each of them is an
\astar search in \emph{space-time}: the state of the drone is a cell together
with the time at which it is there, the drone may wait, and it must avoid the
cells and edges that other drones have reserved. The replanning layer of
\cref{ch:ch24} calls \astar again whenever a local avoidance manoeuvre has
pushed a drone off its nominal route, and the incremental planners of
\cref{ch:ch05} (\lpastar, \dstarlite) and the anytime planner of \cref{ch:ch06}
(\arastar) are best understood as \astar with extra bookkeeping. This chapter is
therefore the foundation of everything that follows in \cref{ch:ch05,ch:ch06,ch:ch08,ch:ch09,ch:ch10}.

The chapter proceeds as follows. \Cref{sec:ch04-intuition} gives the idea,
\cref{sec:ch04-problem} the definitions of admissible and consistent heuristics
and the standard grid heuristics. \Cref{sec:ch04-algorithm} presents the
graph-search version of \astar with a closed set, \cref{sec:ch04-example} runs
it by hand on a $6 \times 6$ grid, and \cref{sec:ch04-properties} proves the
optimality theorems that the Week~1 milestone of the study plan
(\cref{ch:appA}) asks you to be able to explain exactly. \Cref{sec:ch04-variants}
covers tie-breaking, weighted \astar and an experiment on random grids.
\Cref{sec:ch04-spacetime} is a full treatment of space-time \astar, the exact
form of the search used by the multi-agent planners of \cref{ch:ch08,ch:ch09}.
\Cref{sec:ch04-implementation} discusses the implementation and
\cref{sec:ch04-drone} its place in the drone system.

% ---------------------------------------------------------------------
\section{The idea in plain words}
\label{sec:ch04-intuition}

Dijkstra's algorithm always expands the node that is cheapest to reach from the
start. Its frontier grows like a ripple in a pond: a disc around the start that
spreads evenly in all directions until it touches the goal
(\cref{fig:ch04-idea}, left). \astar adds a second ingredient. For every node
$n$ it keeps the cost $\gcost(n)$ of the best path found so far from the start
\emph{and} an estimate $\hcost(n)$ of the cost that remains from $n$ to the goal.
Their sum
\begin{equation}
  \fcost(n) = \gcost(n) + \hcost(n)
  \label{eq:ch04-f}
\end{equation}
estimates the cost of the cheapest solution that passes through $n$. \astar is a
\textbf{best-first search}\index{best-first search} on $\fcost$: it always
expands the node with the smallest $\fcost$.\index{A*!f=g+h@$f=g+h$} The function
$\hcost$ is called a \textbf{heuristic}\index{heuristic} because it is a rule of
thumb, computed without looking at the obstacles, for example the straight-line
distance to the goal. A node far from the goal has a large $\hcost$ and is
expanded late or not at all. The frontier stops being a disc and becomes an
elongated region pointing at the goal (\cref{fig:ch04-idea}, right).

\begin{figure}[tb]
  \centering
  \inputfigure{ch04/idea}
  \caption{Left: Dijkstra expands nodes in order of $\gcost$, so its closed set
  is a disc around the start $s$ that reaches the goal $\gamma$ only when it has
  covered every cheaper node. Right: \astar expands nodes in order of
  $\fcost = \gcost + \hcost$. With the straight-line distance as $\hcost$, the
  nodes with $\fcost \le C^*$ lie inside an ellipse with foci $s$ and $\gamma$
  when the space is obstacle-free, so that $\gcost^*$ is the straight-line
  distance; the small obstacle drawn here deforms the region but does not
  change the picture. The search never leaves it. For the node $n$, the solid segment is the known
  cost $\gcost(n)$ and the dashed segment the estimate $\hcost(n)$.}
  \label{fig:ch04-idea}
\end{figure}

Two questions decide whether this idea is safe. First, can a heuristic mislead
the search into returning a path that is not the shortest? Yes, if it
\emph{overestimates}: a node on the true shortest path may then look worse than
it is and never be expanded. Second, can a node be expanded twice? Yes, if the
heuristic is not \emph{consistent}, because a cheaper path to an already
expanded node may show up later. The two conditions that rule out these
problems, admissibility and consistency, are the heart of this chapter.

\begin{keyidea}
Expand the node with the smallest $\fcost = \gcost + \hcost$. If $\hcost$ never
overestimates the remaining cost, the first goal that \astar \emph{pops} is
reached by an optimal path, provided a closed node may be re-opened when a
cheaper path to it turns up (\cref{alg:ch04-astar},
line~\ref{alg:ch04-astar:reopen}). If in addition $\hcost$ obeys the triangle inequality along
every edge, no node is ever expanded twice, and \astar is Dijkstra's algorithm
with its effort focused towards the goal.
\end{keyidea}

% ---------------------------------------------------------------------
\section{Problem statement and notation}
\label{sec:ch04-problem}

\subsection{Search problems and heuristics}

\begin{definition}[Heuristic search problem]
\label{def:ch04-problem}
A \textbf{heuristic search problem}\index{search problem!heuristic} consists of a
finite directed graph $G = (V, E)$ with edge costs $c(u, v) \ge 0$, a start node
$s \in V$, a non-empty set of goal nodes $V_\gamma \subseteq V$ (usually a single
node $\gamma$), and a \textbf{heuristic function} $\hcost \colon V \to
\R_{\ge 0}$. The cost of a path is the sum of its edge costs. We write
$\gcost^*(n)$ for the cost of a cheapest path from $s$ to $n$,
$\hcost^* \colon V \to \R_{\ge 0} \cup \{\infty\}$ for the cost of a cheapest
path from $n$ to a goal ($\hcost^*(n) = \infty$ if no goal is reachable; the
heuristic $\hcost$ itself is assumed finite everywhere), and $C^* = \hcost^*(s)$ for
the optimal solution cost.\index{$h^*$, true cost to go}\index{$C^*$, optimal cost}
\end{definition}

During the search, $\gcost(n)$ denotes the cost of the best path from $s$ to
$n$ found \emph{so far}; it satisfies $\gcost(n) \ge \gcost^*(n)$, with equality
once the best path has been found. The successor function $\Succ(n)$ returns the
pairs $(n', c(n, n'))$ for all edges $(n, n') \in E$, exactly as in
\cref{ch:ch03}. For grids the nodes are cells $(x, y)$, $x$ the column and
$y$ the row, with $(0, 0)$ in the bottom-left corner as in \cref{ch:ch02}.

\begin{definition}[Admissible heuristic]
\label{def:ch04-admissible}
A heuristic $\hcost$ is \textbf{admissible}\index{heuristic!admissible} if it never
overestimates the true cost to go:
\begin{equation}
  0 \le \hcost(n) \le \hcost^*(n) \quad \text{for every } n \in V.
  \label{eq:ch04-admissible}
\end{equation}
In particular $\hcost(\gamma) = 0$ for every goal node $\gamma$.
\end{definition}

\begin{definition}[Consistent heuristic]
\label{def:ch04-consistent}
A heuristic $\hcost$ is \textbf{consistent}\index{heuristic!consistent} (also
called \textbf{monotone}\index{heuristic!monotone}) if $\hcost(\gamma) = 0$ for
every goal node $\gamma$ and, for every edge $(n, n') \in E$,
\begin{equation}
  \hcost(n) \le c(n, n') + \hcost(n').
  \label{eq:ch04-consistent}
\end{equation}
\end{definition}

\Cref{eq:ch04-consistent} is a triangle inequality: the estimate from $n$ may not
exceed the cost of first stepping to $n'$ and estimating from there. An
equivalent formulation is that $\fcost$ never decreases along an edge when
$\gcost$ is computed along the path: $\fcost(n') = \gcost(n) + c(n, n') +
\hcost(n') \ge \gcost(n) + \hcost(n) = \fcost(n)$. Consistency is the stronger
property.

\begin{lemma}[Consistency implies admissibility]
\label{thm:ch04-consistent-admissible}
Every consistent heuristic is admissible. The converse is false.
\end{lemma}
\begin{proof}
Let $\hcost$ be consistent and let $n \in V$. If no goal is reachable from $n$,
then $\hcost^*(n) = \infty$ and \cref{eq:ch04-admissible} holds. Otherwise let
$n = n_0, n_1, \dots, n_k = \gamma$ be a cheapest path from $n$ to a goal, so that
$\hcost^*(n) = \sum_{i=0}^{k-1} c(n_i, n_{i+1})$. Applying
\cref{eq:ch04-consistent} to each edge in turn,
\[
  \hcost(n_0) \le c(n_0, n_1) + \hcost(n_1)
             \le c(n_0, n_1) + c(n_1, n_2) + \hcost(n_2)
             \le \dots
             \le \sum_{i=0}^{k-1} c(n_i, n_{i+1}) + \hcost(\gamma) = \hcost^*(n),
\]
because $\hcost(\gamma) = 0$. For the converse, take the graph with nodes
$s, a, b, \gamma$ and edges $s \to a$ (cost 1), $s \to b$ (cost 3), $a \to b$
(cost 1), $b \to \gamma$ (cost 2). The heuristic $\hcost(a) = 3$,
$\hcost(s) = \hcost(b) = \hcost(\gamma) = 0$ is admissible, since
$\hcost^*(a) = 3$, but $\hcost(a) = 3 > c(a, b) + \hcost(b) = 1$ violates
\cref{eq:ch04-consistent}. We return to this graph in
\cref{sec:ch04-properties}.
\end{proof}

\subsection{Heuristics for grids}
\label{sec:ch04-grid-heuristics}

Let $n = (x_n, y_n)$ be a cell and $\gamma = (x_\gamma, y_\gamma)$ the goal, and
write $d_x = \abs{x_n - x_\gamma}$ and $d_y = \abs{y_n - y_\gamma}$. Four
distances are in everyday use (\cref{fig:ch04-heuristics}):
\begin{align}
  \hcost_{\mathrm{M}}(n) &= d_x + d_y
    && \text{Manhattan distance,}\index{heuristic!Manhattan} \label{eq:ch04-manhattan}\\
  \hcost_{\mathrm{O}}(n) &= \max(d_x, d_y) + (\sqrt{2} - 1)\min(d_x, d_y)
    && \text{octile distance,}\index{heuristic!octile} \label{eq:ch04-octile}\\
  \hcost_{\mathrm{E}}(n) &= \sqrt{d_x^2 + d_y^2}
    && \text{Euclidean distance,}\index{heuristic!Euclidean} \label{eq:ch04-euclid}\\
  \hcost_{\mathrm{C}}(n) &= \max(d_x, d_y)
    && \text{Chebyshev distance.}\index{heuristic!Chebyshev} \label{eq:ch04-chebyshev}
\end{align}
The Manhattan distance is the length of a cheapest path on an empty
4-connected grid, where every move costs~1. The octile distance is the length
of a cheapest path on an empty 8-connected grid, where a straight move costs~1
and a diagonal move $\sqrt{2}$: use $\min(d_x, d_y)$ diagonal moves and then
$\max(d_x, d_y) - \min(d_x, d_y)$ straight moves. The Chebyshev distance is the
same count when diagonal moves also cost~1. On an 8-connected grid we use the
successor rule of \cref{def:ch02-grid} throughout: a diagonal move is allowed
only when both orthogonal neighbours it passes are free, so a drone of positive
size never clips an obstacle corner. Forbidding those moves only lengthens
paths, so it does not endanger the admissibility of any heuristic below, and
\code{grid\_successors()} in \code{ch04\_astar.py} enforces it in every search
of this chapter, including the space-time search of
\cref{sec:ch04-spacetime}. Each of the four is the norm of the
displacement vector ($\ell_1$, octagonal, $\ell_2$ and $\ell_\infty$ norm), so
each satisfies the triangle inequality.

\begin{figure}[tb]
  \centering
  \inputfigure{ch04/heuristics}
  \caption{The four grid heuristics from the cell $n$ to the goal $\gamma$ with
  $d_x = 5$ and $d_y = 3$. The Manhattan path (orange, dashed) has length~8; the
  octile path (blue) uses three diagonal moves and two straight ones, length
  $5 + 3(\sqrt{2} - 1) \approx 6.24$; the Euclidean distance (green) is
  $\sqrt{34} \approx 5.83$; the Chebyshev distance counts only the longer side,
  $5$. On a 4-connected grid all four are admissible; on an 8-connected grid the
  Manhattan distance overestimates every diagonal move.}
  \label{fig:ch04-heuristics}
\end{figure}

\begin{proposition}[Metric heuristics are consistent]
\label{thm:ch04-metric}
Let $d$ be a metric on $V$ such that $d(n, n') \le c(n, n')$ for every edge
$(n, n') \in E$, and let $\hcost(n) = d(n, \gamma)$. Then $\hcost$ is
consistent, hence admissible.
\end{proposition}
\begin{proof}
$\hcost(\gamma) = d(\gamma, \gamma) = 0$, and by the triangle inequality
$\hcost(n) = d(n, \gamma) \le d(n, n') + d(n', \gamma) \le c(n, n') + \hcost(n')$.
\end{proof}

To apply \cref{thm:ch04-metric} we only have to compare each distance with the
cost of a single move. A straight move costs~1, and all four distances assign
$1$ to it. A diagonal move on an 8-connected grid costs $\sqrt{2}$: the octile
and Euclidean distances assign $\sqrt{2}$, the Chebyshev distance $1$, but the
Manhattan distance assigns $2 > \sqrt{2}$. \Cref{tab:ch04-heuristics} summarises
the result. The Manhattan distance is the classic example of a heuristic that is
perfect for one grid and wrong for another.

\begin{table}[tb]
  \centering\footnotesize
  \caption{Grid heuristics and their properties. ``Exact'' means equal to
  $\hcost^*$ on an obstacle-free grid. A heuristic that is exact on empty grids
  stays admissible with obstacles, because obstacles can only lengthen paths.}
  \label{tab:ch04-heuristics}
  \begin{tabular}{@{}lll@{}}
  \toprule
  Heuristic & 4-connected (cost 1) & 8-connected (costs $1$, $\sqrt{2}$)\\
  \midrule
  Manhattan $d_x + d_y$ & exact on empty grids; consistent & \textbf{not admissible} (diagonal: $2 > \sqrt{2}$)\\
  Octile $\max + (\sqrt{2}-1)\min$ & not exact; consistent (below Manhattan) & exact on empty grids; consistent\\
  Euclidean $\sqrt{d_x^2 + d_y^2}$ & consistent & consistent; dominated by octile\\
  Chebyshev $\max(d_x, d_y)$ & consistent; weakest of the four & consistent; exact if diagonals cost 1\\
  \bottomrule
  \end{tabular}
\end{table}

\begin{pitfall}[The Manhattan distance on an 8-connected grid]
This is the most common inadmissible heuristic in robotics code, and it is
invisible: the search runs, returns a path, and the path is simply too long.
The reason is one line of \cref{tab:ch04-heuristics}: a diagonal move costs
$\sqrt{2} \approx 1.41$ but changes $d_x + d_y$ by~$2$, so
$\hcost_{\mathrm{M}}(n) > c(n,n') + \hcost_{\mathrm{M}}(n')$ and, worse,
$\hcost_{\mathrm{M}}$ can exceed $\hcost^*$. \Cref{thm:ch04-optimal} then does
not apply. In the self-test of \code{ch04\_astar.py}, \astar with the Manhattan
heuristic on 8-connected random grids returns a strictly longer path on 6 of 40
instances --- one in seven, silently. The bug is easy to introduce when a
4-connected planner is extended to diagonal moves and nobody revisits the
heuristic. Use the octile distance \cref{eq:ch04-octile} on 8-connected maps,
and test optimality against a Dijkstra run on random instances, as the self-test
does.
\end{pitfall}

None of the four sees obstacles, so with walls in the way they all
underestimate, sometimes badly: the worked example of \cref{sec:ch04-example}
shows the search walking into a dead end because the Manhattan distance does
not know about a wall. The exact heuristic $\hcost^*$ itself can be computed by a
backward Dijkstra search from the goal, as in \cref{ch:ch03}. That costs a full
search and is pointless for a single query, but it becomes the heuristic of
choice in space-time (\cref{sec:ch04-spacetime}), where the same static
distances serve every time layer and every constrained re-run of the search.

% ---------------------------------------------------------------------
\section{The algorithm}
\label{sec:ch04-algorithm}

\Cref{alg:ch04-astar} is \astar as a \emph{graph search}: it remembers the nodes
it has expanded in the set \Closed\index{closed set} so that a node is not
expanded again unless a strictly cheaper path to it is found. The nodes that
have been generated but not yet expanded form the frontier \Open, a priority
queue keyed on $\fcost$.\index{open set} The queue is \emph{lazy}: when the
$\gcost$ of a node improves, the algorithm pushes a fresh entry instead of
searching the queue for the old one, and the old entry is discarded when it is
popped and found to be stale.\index{priority queue!lazy deletion} The weight $w$
is $1$ for plain \astar; \cref{sec:ch04-variants} discusses $w > 1$.

\begin{algorithm}[htb]
\caption{\astar graph search with a closed set and a lazy priority queue}
\label{alg:ch04-astar}
\KwIn{graph $G = (V, E)$ with costs $c$, start $s$, goal test, heuristic $\hcost$, weight $w \ge 1$ ($w = 1$ for \astar)}
\KwOut{a cheapest path from $s$ to a goal, or \emph{failure}}
\Fn{\AstarSearch{$G$, $s$, $\mathit{isGoal}$, $\hcost$, $w$}}{
  $\gcost(s) \gets 0$; parent$(s) \gets$ \KwNil; $\gcost(n) \gets \infty$ for all other nodes (implicitly)\;
  \Open $\gets$ priority queue holding $s$ with key $(w\,\hcost(s),\, 0)$; \Closed $\gets \emptyset$\label{alg:ch04-astar:init}\;
  \While{\Open $\ne \emptyset$}{
    $(n, \gcost_n) \gets$ pop the entry with the smallest key $(\fcost, -\gcost)$ from \Open\label{alg:ch04-astar:pop}\;
    \If{$\gcost_n > \gcost(n)$}{\KwContinue\tcp*{stale entry: a cheaper path to $n$ was found later}\label{alg:ch04-astar:stale}}
    \If{$n$ is a goal}{\KwRet \AstarPath{parent, $n$}\label{alg:ch04-astar:goal}\;}
    add $n$ to \Closed\label{alg:ch04-astar:close}\;
    \ForEach{$(n', c(n, n')) \in \Succ(n)$}{
    $\gcost' \gets \gcost(n) + c(n, n')$\;
      \If{$\gcost' < \gcost(n')$}{\label{alg:ch04-astar:relax}
        \If{$n' \in$ \Closed}{remove $n'$ from \Closed\tcp*{re-open $n'$; never happens if $\hcost$ is consistent}\label{alg:ch04-astar:reopen}}
        $\gcost(n') \gets \gcost'$; parent$(n') \gets n$\;
        push $(n', \gcost')$ into \Open with key $(\gcost' + w\,\hcost(n'),\, -\gcost')$\label{alg:ch04-astar:push}\;
      }
    }
  }
  \KwRet \emph{failure}\;
}
\end{algorithm}

\paragraph{Walkthrough.}
Line~\ref{alg:ch04-astar:init} seeds the frontier with the start, whose
$\fcost$ is $\hcost(s)$. The loop repeatedly takes the most promising entry
(line~\ref{alg:ch04-astar:pop}). The key is the pair $(\fcost, -\gcost)$, compared
lexicographically: smaller $\fcost$ first and, among equal $\fcost$, larger
$\gcost$ first. This tie-breaking rule\index{tie-breaking} is not needed for
correctness, but \cref{sec:ch04-variants} shows that it can make a tenfold
difference in effort. Line~\ref{alg:ch04-astar:stale} implements lazy deletion:
a node may have several entries in \Open, one per improvement of its $\gcost$;
only the entry carrying the current value is valid. The goal test happens when
the goal is \emph{popped} (line~\ref{alg:ch04-astar:goal}), not when it is first
generated. This detail is what makes the optimality proof of
\cref{sec:ch04-properties} work; testing at generation time returns the first
path found to the goal, which need not be the cheapest. Line~\ref{alg:ch04-astar:close}
records the expansion. The inner loop relaxes every outgoing edge
(line~\ref{alg:ch04-astar:relax}) exactly as Dijkstra's algorithm does; the
only differences from \cref{ch:ch03} are the key of the priority queue and
line~\ref{alg:ch04-astar:reopen}, which re-opens a closed node when a cheaper
path to it appears. With a consistent heuristic this line is never executed
(\cref{thm:ch04-consistent}), and many implementations drop it together with
the ability to correct a closed node; with an admissible but inconsistent
heuristic it is required for optimality (\cref{thm:ch04-optimal}). The
function \AstarPath follows the parent pointers back from the goal to $s$ and
reverses the list.

Setting $\hcost \equiv 0$ turns the key into $(\gcost, -\gcost)$ and
\cref{alg:ch04-astar} into Dijkstra's algorithm with early exit; every property
of \cref{ch:ch03} is the special case $\hcost = 0$ of a property in this
chapter.\index{A*!as generalisation of Dijkstra}

% ---------------------------------------------------------------------
\section{A worked example}
\label{sec:ch04-example}

\begin{example}[\astar on a $6 \times 6$ grid]
\label{ex:ch04-grid}
\Cref{fig:ch04-example} shows a $6 \times 6$ 4-connected grid with two walls:
the cells $(2, 0)$ to $(2, 3)$ and $(4, 1)$ to $(4, 4)$ are blocked. The start
is $s = (0, 0)$ in the bottom-left corner and the goal $\gamma = (5, 2)$. Every
move costs~1 and the heuristic is the Manhattan distance
\cref{eq:ch04-manhattan}, so $\hcost(s) = 5 + 2 = 7$. Ties in $\fcost$ are
broken towards the larger $\gcost$, and among equal keys the entry pushed first
is popped first. The instance is \code{worked\_example\_grid()} in
\code{ch04\_astar.py}, and every number below is produced by its self-test.
\end{example}

\begin{figure}[tb]
  \centering
  \inputfigure{ch04/example}
  \caption{\Cref{ex:ch04-grid} after termination. Grey cells are \Closed
  (expanded), blue cells are still in \Open; the bold number is $\fcost$ at the
  time of expansion (for Dijkstra $\fcost = \gcost$) and the small grey number
  in the corner is the expansion order. The blue line is the returned path of
  cost~13. (a) \astar with the Manhattan heuristic expands 21 cells; note the
  three wasted expansions $(3, 3)$, $(3, 2)$, $(3, 1)$ below the junction at
  $(3, 4)$, where the heuristic cannot see the wall at $x = 4$. (b) Dijkstra's
  algorithm ($\hcost = 0$) expands 26 cells, five more than \astar: the top-row
  cells $(0, 5)$, $(1, 5)$, $(2, 5)$ and the bottom-row cells $(3, 0)$,
  $(4, 0)$, which \astar either leaves in \Open or never generates.}
  \label{fig:ch04-example}
\end{figure}

\Cref{tab:ch04-trace} lists the first ten expansions. Step~1 expands the start
$(0, 0)$; steps~2--4 then walk up the column $x = 1$: every cell there has $\fcost = 7$, because moving towards the
goal decreases $\hcost$ by exactly the move cost. At $(1, 2)$ the search would
like to continue east, but $(2, 2)$ is blocked; going north to $(1, 3)$ moves
away from the goal row, and $\fcost$ rises to $9$. \astar therefore turns to the
remaining cells with $\fcost = 7$, the column $x = 0$ (steps~5 and~6), before it
accepts any cell with $\fcost = 9$ (steps~7 and~8). From step~9 on the
frontier has $\fcost = 11$. The search passes the first wall at $(2, 4)$ and,
led by the decreasing $\hcost$, dives down the column $x = 3$ as far as
$(3, 1)$, where the second wall stops it: the three cells $(3, 3), (3, 2),
(3, 1)$ below the junction at $(3, 4)$ are a dead end that the Manhattan
distance cannot foresee. The cell $(0,4)$ carries $\fcost = 11$ and $(3,1)$
carries $\fcost = 13$, so $(0,4)$ is popped first and $(3,1)$ one pop later; the
expansion order is $\dots, (3, 4), (3, 3), (3, 2), (0, 4), (3, 1), (3, 5),
\dots$. From $(3,1)$ on the frontier is at $\fcost = 13$: the search crosses the
top row $(3, 5), (4, 5), (5, 5)$ and descends to the goal, which is popped at
step~21 with $\gcost = 13$. At that moment three cells with
$\fcost = 13$ but smaller $\gcost$ are still in \Open together with $(3, 0)$ at
$\fcost = 15$; none of them is ever expanded. Dijkstra's algorithm needs 26
expansions for the same instance, and \astar with ties broken towards the
\emph{smaller} $\gcost$ needs 24. On the 8-connected version of the grid, \astar
with the octile heuristic finds a path of cost $9 + 2\sqrt{2} \approx 11.83$
with 24 expansions.

\begin{table}[tb]
  \centering\small
  \caption{The first ten expansions of \cref{alg:ch04-astar} on
  \cref{ex:ch04-grid}. The last column lists the cells in \Open after the step
  with their current $\fcost$; a cell whose $\gcost$ improves, such as $(0, 2)$
  in step~5, keeps a single valid entry.}
  \label{tab:ch04-trace}
  \begin{tabular}{@{}rlrrrl@{}}
  \toprule
  Step & Expanded $n$ & $\gcost$ & $\hcost$ & $\fcost$ & \Open after the step (cell$:\fcost$)\\
  \midrule
   1 & $(0,0)$ & 0 & 7 &  7 & $(0,1){:}7$, $(1,0){:}7$\\
   2 & $(1,0)$ & 1 & 6 &  7 & $(0,1){:}7$, $(1,1){:}7$\\
   3 & $(1,1)$ & 2 & 5 &  7 & $(0,1){:}7$, $(1,2){:}7$\\
   4 & $(1,2)$ & 3 & 4 &  7 & $(0,1){:}7$, $(0,2){:}9$, $(1,3){:}9$\\
   5 & $(0,1)$ & 1 & 6 &  7 & $(0,2){:}7$, $(1,3){:}9$\\
   6 & $(0,2)$ & 2 & 5 &  7 & $(0,3){:}9$, $(1,3){:}9$\\
   7 & $(1,3)$ & 4 & 5 &  9 & $(0,3){:}9$, $(1,4){:}11$\\
   8 & $(0,3)$ & 3 & 6 &  9 & $(0,4){:}11$, $(1,4){:}11$\\
   9 & $(1,4)$ & 5 & 6 & 11 & $(0,4){:}11$, $(1,5){:}13$, $(2,4){:}11$\\
  10 & $(2,4)$ & 6 & 5 & 11 & $(0,4){:}11$, $(1,5){:}13$, $(2,5){:}13$, $(3,4){:}11$\\
  \bottomrule
  \end{tabular}
\end{table}

Two things in the table deserve a second look. In step~4 the cell $(0, 2)$ enters
\Open with $\gcost = 3$ (through $(1, 2)$) and $\fcost = 9$; in step~5 the
expansion of $(0, 1)$ finds the cheaper route with $\gcost = 2$, the entry is
replaced and $\fcost$ drops to $7$. This is the relaxation of
line~\ref{alg:ch04-astar:relax}, and it happens while $(0, 2)$ is still open,
so no re-opening is involved. Second, the $\fcost$ of the expanded nodes reads
$7, 7, 7, 7, 7, 7, 9, 9, 11, 11, \dots, 13$: it never decreases. This is not
luck; \cref{thm:ch04-consistent} shows that it holds for every consistent
heuristic.

% ---------------------------------------------------------------------
\section{Properties: correctness, optimality, complexity}
\label{sec:ch04-properties}

Throughout this section $G$ is finite, costs are non-negative and
\cref{alg:ch04-astar} runs with $w = 1$. The proofs follow Hart, Nilsson and
Raphael~\cite{hart1968formal,hart1972correction} in the form given by Russell
and Norvig~\cite{russell2020aima}; Pearl~\cite{pearl1984heuristics} treats the
general case.

\subsection{Termination and completeness}

\begin{proposition}[Termination and completeness]
\label{thm:ch04-terminates}
\Cref{alg:ch04-astar} terminates. It returns a path if and only if some goal is
reachable from $s$.
\end{proposition}
\begin{proof}
Every push into \Open strictly decreases $\gcost(n')$ for some node $n'$
(line~\ref{alg:ch04-astar:relax}). The path recorded by the parent pointers of
$n'$ is simple: if it returned to $n'$, the non-negative costs along the cycle
would give $\gcost(n) + c(n, n') \ge \gcost(n')$, contradicting the strict
decrease. A finite graph has finitely many simple paths, so each node's
$\gcost$ can decrease only finitely often, there are finitely many pushes, and
the loop ends. If no goal is reachable, no goal is ever popped and the
algorithm returns \emph{failure} when \Open runs empty. If a goal is reachable,
\cref{thm:ch04-invariant} below shows that \Open is non-empty until a goal is
expanded, so the algorithm returns a path.
\end{proof}

\subsection{Optimality with an admissible heuristic}

The proofs rest on one invariant. It concerns a cheapest path to an arbitrary
node $m$; we apply it with $m$ a goal and later with $m$ an ordinary node.

\begin{lemma}[An optimal-path node is always open]
\label{thm:ch04-invariant}
Let $P = (s = n_0, n_1, \dots, n_k = m)$ be a cheapest path from $s$ to a node
$m$. At any time before \cref{alg:ch04-astar} expands $m$ with
$\gcost(m) = \gcost^*(m)$, \Open contains a node $n_i$ of $P$ with
$\gcost(n_i) = \gcost^*(n_i)$.
\end{lemma}
\begin{proof}
Every prefix of a cheapest path is a cheapest path, so
$\gcost^*(n_{i+1}) = \gcost^*(n_i) + c(n_i, n_{i+1})$. Let $i$ be the smallest
index such that $n_i$ has \emph{not} yet been expanded while holding
$\gcost(n_i) = \gcost^*(n_i)$; the index exists because $m = n_k$ qualifies. If
$i = 0$, then $s$ has not been expanded (it always holds $\gcost(s) = 0 =
\gcost^*(s)$), so its initial entry is still in \Open. If $i > 0$, then
$n_{i-1}$ was expanded with the optimal value, and at that moment
line~\ref{alg:ch04-astar:relax} set $\gcost(n_i) \le \gcost^*(n_{i-1}) +
c(n_{i-1}, n_i) = \gcost^*(n_i)$. Since every $\gcost$ value is the cost of an
actual path, $\gcost(n_i) \ge \gcost^*(n_i)$, hence $\gcost(n_i) =
\gcost^*(n_i)$ from then on. The push that established this value (either at
that relaxation, which also re-opened $n_i$ if it was closed, or earlier) put an
entry with $\gcost^*(n_i)$ into \Open. That entry is never stale, and by the
choice of $i$ it has not been popped for expansion; so it is still in \Open.
\end{proof}

\begin{theorem}[Optimality of \astar; Hart, Nilsson and Raphael 1968]
\label{thm:ch04-optimal}
If $\hcost$ is admissible and a goal is reachable from $s$, then
\cref{alg:ch04-astar} (with re-opening, line~\ref{alg:ch04-astar:reopen})
returns a path of cost $C^*$.\index{A*!optimality theorem}
\end{theorem}
\begin{proof}
By \cref{thm:ch04-terminates} the algorithm returns a path to some goal
$\gamma'$; its cost is $\gcost(\gamma')$ at the moment $\gamma'$ is popped.
Suppose $\gcost(\gamma') > C^*$. Let $P = (s, \dots, \gamma)$ be an optimal
solution, of cost $C^*$. The goal $\gamma$ has not been expanded with
$\gcost(\gamma) = C^*$, since the algorithm would have stopped there. By
\cref{thm:ch04-invariant}, at the moment $\gamma'$ is popped \Open contains a
node $n_i$ of $P$ with $\gcost(n_i) = \gcost^*(n_i)$. The rest of $P$ from $n_i$
to $\gamma$ costs $C^* - \gcost^*(n_i)$ and is one particular path to a goal,
so $\hcost^*(n_i) \le C^* - \gcost^*(n_i)$, and admissibility gives
\[
  \fcost(n_i) = \gcost^*(n_i) + \hcost(n_i)
             \le \gcost^*(n_i) + \hcost^*(n_i) \le C^*.
\]
But $\fcost(\gamma') = \gcost(\gamma') + \hcost(\gamma') = \gcost(\gamma') >
C^* \ge \fcost(n_i)$, so line~\ref{alg:ch04-astar:pop} would have popped
$n_i$, or some other entry with a smaller key, before $\gamma'$. This
contradiction shows $\gcost(\gamma') \le C^*$, and $\gcost(\gamma') \ge C^*$
because it is the cost of a real path.
\end{proof}

The proof used admissibility only in one place: to bound $\fcost(n_i)$ by
$C^*$. It used re-opening, through \cref{thm:ch04-invariant}, to make sure
that the optimal-path node is in \Open even if it had been expanded earlier
with a worse $\gcost$. With a consistent heuristic the second precaution is
unnecessary.

\subsection{Consistent heuristics: no node is expanded twice}

\begin{theorem}[Consistent heuristics]
\label{thm:ch04-consistent}
Let $\hcost$ be consistent. Then in \cref{alg:ch04-astar}:
\begin{enumerate}[label=(\roman*)]
  \item whenever a node $n$ is expanded, $\gcost(n) = \gcost^*(n)$; hence
        line~\ref{alg:ch04-astar:reopen} is never executed and every node is
        expanded at most once;
  \item the $\fcost$ values of the expanded nodes, in order of expansion, are
        non-decreasing.
\end{enumerate}
\index{A*!consistency}\index{re-opening}
\end{theorem}
\begin{proof}
(i) Suppose $n$ is popped for expansion with $\gcost(n) > \gcost^*(n)$. Let
$P = (s, \dots, n)$ be a cheapest path to $n$. By \cref{thm:ch04-invariant}
with $m = n$, \Open contains a node $n_i$ of $P$ with $\gcost(n_i) =
\gcost^*(n_i)$, and $n_i \ne n$ because $\gcost(n)$ is not optimal. Applying
\cref{eq:ch04-consistent} along the edges of $P$ from $n_i$ to $n$, exactly as
in the proof of \cref{thm:ch04-consistent-admissible}, gives $\hcost(n_i) \le
\bigl(\gcost^*(n) - \gcost^*(n_i)\bigr) + \hcost(n)$. Therefore
\[
  \fcost(n_i) = \gcost^*(n_i) + \hcost(n_i) \le \gcost^*(n) + \hcost(n)
             < \gcost(n) + \hcost(n) = \fcost(n),
\]
and line~\ref{alg:ch04-astar:pop} would have chosen $n_i$ instead of $n$: a
contradiction. So every expanded node carries its optimal $\gcost$, no later
path can be cheaper, and a closed node is never re-opened.

(ii) Let $n$ be expanded and let $n'$ be the node expanded next. When $n$ was
popped, every other entry in \Open had $\fcost \ge \fcost(n)$. The expansion of
$n$ changes \Open in two ways: it pushes successors $n''$ with
$\fcost(n'') = \gcost(n) + c(n, n'') + \hcost(n'') \ge \gcost(n) + \hcost(n) =
\fcost(n)$ by consistency, and it leaves the other entries unchanged. Hence
every entry in \Open has $\fcost \ge \fcost(n)$ and in particular
$\fcost(n') \ge \fcost(n)$.
\end{proof}

\begin{corollary}[Which nodes \astar expands]
\label{thm:ch04-surely}
With a consistent heuristic, \astar expands every node $n$ with
$\gcost^*(n) + \hcost(n) < C^*$, expands no node with $\gcost^*(n) + \hcost(n)
> C^*$, and may or may not expand nodes with $\gcost^*(n) + \hcost(n) = C^*$
depending on tie-breaking.
\end{corollary}
\begin{proof}
By \cref{thm:ch04-consistent}(ii) and $\fcost(\gamma) = C^*$ at termination, no
expanded node has $\fcost > C^*$. Conversely let $\gcost^*(n) + \hcost(n) <
C^*$ and let $P$ be a cheapest path to $n$. Along $P$ the value
$\gcost^* + \hcost$ is non-decreasing (consistency), so it is below $C^*$ at
every node of $P$. By induction along $P$, each node of $P$ is generated with
its optimal $\gcost$ by the expansion of its predecessor and, having
$\fcost < C^* = \fcost(\gamma)$, is popped before the goal.
\end{proof}

\Cref{thm:ch04-consistent}(i) is the reason why the closed set can be a plain
``expanded'' flag when $\hcost$ is consistent, and part~(ii) why the trace in
\cref{tab:ch04-trace} shows non-decreasing $\fcost$;
\cref{thm:ch04-surely} is why the $\fcost$ values in \cref{fig:ch04-example}
never exceed $13$.

\subsection{What goes wrong without the conditions}

\begin{figure}[tb]
  \centering
  \inputfigure{ch04/inadmissible}
  \caption{An inadmissible heuristic misleads \astar. The value $\hcost(a) = 4$
  overestimates $\hcost^*(a) = 1$, so $a$ looks worse than the detour through
  $b$ and is never expanded (blue fill: still in \Open). \astar returns the path
  of cost~3 although a path of cost~2 exists.}
  \label{fig:ch04-inadmissible}
\end{figure}

\paragraph{Inadmissible heuristics.}
\Cref{fig:ch04-inadmissible} shows the smallest kind of failure. The heuristic
is wrong at a single node, $a$, where it says $4$ while the true cost to go is
$1$. The search expands $s$, then $b$ (with $\fcost = 2$), then the goal with
$\gcost = 3$, and it stops. The node $a$ sits in \Open with $\fcost = 5$ and
the optimal path $s \to a \to \gamma$ of cost~2 is never found. Nothing in the
run signals the error: an overestimate makes \astar \emph{confidently} wrong,
which is why the Manhattan heuristic on an 8-connected grid
(\cref{tab:ch04-heuristics}) is a real bug and not a matter of taste.
\index{heuristic!inadmissible}

\paragraph{Admissible but inconsistent heuristics.}
Take the graph from the proof of \cref{thm:ch04-consistent-admissible}:
$s \to a$ (1), $s \to b$ (3), $a \to b$ (1), $b \to \gamma$ (2), with
$\hcost(a) = 3$ and $\hcost = 0$ elsewhere. \astar expands $s$
($\fcost = 0$) and pushes $a$ with $\fcost = 1 + 3 = 4$ and $b$ with
$\fcost = 3 + 0 = 3$. It expands $b$ first and pushes $\gamma$ with
$\gcost = 5$. Next it expands $a$ and finds $b$ with $\gcost = 2 < 3$: $b$ is
closed, and line~\ref{alg:ch04-astar:reopen} re-opens it. The expansion of $b$
(now $\fcost = 2$) improves $\gamma$ to $\gcost = 4$, and the goal is
returned with the optimal cost~4. The sequence of $\fcost$ values,
$0, 3, 4, 2, 4$, is not monotone, and $b$ is expanded twice. If re-opening is
disabled, the same run returns the path $s \to b \to \gamma$ of cost~5. The
self-test of \code{ch04\_astar.py} reproduces both runs.

\subsection{Complexity}

With a consistent heuristic every node is expanded at most once and every
edge relaxed at most once, so \cref{alg:ch04-astar} performs at most $\abs{E}$
pushes and $\abs{V}$ expansions; with a binary heap the running time is
$\bigO{(\abs{V} + \abs{E}) \log \abs{V}}$ and the memory $\bigO{\abs{V} +
\abs{E}}$, the same bounds as for Dijkstra's algorithm in \cref{ch:ch03}. The
bounds are worst cases. What decides the actual effort is the number of
expanded nodes, that is, by \cref{thm:ch04-surely}, the number of nodes with
$\gcost^*(n) + \hcost(n) < C^*$. With $\hcost = 0$ this is every node cheaper
than the goal; with the exact heuristic $\hcost = \hcost^*$ it is empty, and
with larger-$\gcost$-first tie-breaking \astar then expands only the nodes of
one optimal path. Between these extremes the effort grows with the error
$\hcost^* - \hcost$: on a grid with obstacles every detour that the heuristic
does not foresee adds a band of cells with $\fcost < C^*$.
On an \emph{implicit} graph, one whose nodes are generated on demand and whose
size is exponential in the solution depth, the $\abs{V}$ in the bound is that
exponential number, and the heuristic is what decides whether the search is
practical. Pearl~\cite[ch.~6]{pearl1984heuristics} makes this precise: \astar
still expands exponentially many nodes unless the error of the heuristic grows
more slowly than the true cost, roughly $\hcost^*(n) - \hcost(n) =
\bigO{\log \hcost^*(n)}$; a constant \emph{relative} error is not enough, since
a heuristic with $\hcost = (1-\epsilon)\hcost^*$ still gives exponentially many
expansions --- which is why \cref{sec:ch04-variants} buys speed by giving up
optimality rather than by sharpening $\hcost$.
All of this counts \emph{distinct} nodes. If $\hcost$ is admissible but not
consistent, line~\ref{alg:ch04-astar:reopen} may fire again and again, and
Martelli~\cite{martelli1977complexity} builds a family of graphs on which
\astar re-expands nodes exponentially often in the number of nodes, so the
$\bigO{\abs{V}}$ expansion count above holds only under
\cref{thm:ch04-consistent}(i); Pearl~\cite[ch.~3]{pearl1984heuristics}
discusses the algorithms that repair this.
The grids of this book are explicit and small enough that this does not bite,
but the space-time graphs of \cref{sec:ch04-spacetime} and the constraint
trees of \cref{ch:ch09} are where it starts to.
\index{A*!complexity}

\subsection{Heuristic dominance}

\begin{definition}[Dominance]
\label{def:ch04-dominance}
A heuristic $\hcost_2$ \textbf{dominates}\index{heuristic!dominance} $\hcost_1$ if
$\hcost_2(n) \ge \hcost_1(n)$ for every node $n$. If both are admissible,
$\hcost_2$ is the \emph{more informed} of the two.
\end{definition}

\begin{theorem}[More informed, fewer expansions]
\label{thm:ch04-dominance}
Let $\hcost_1$ and $\hcost_2$ be consistent and let $\hcost_2$ dominate
$\hcost_1$. Every node that \astar with $\hcost_2$ expands with
$\fcost < C^*$ is also expanded by \astar with $\hcost_1$. Consequently the
number of expansions with $\hcost_2$ is at most the number with $\hcost_1$, up
to nodes with $\fcost = C^*$.
\end{theorem}
\begin{proof}[Proof sketch]
By \cref{thm:ch04-surely}, \astar with $\hcost_2$ expands with $\fcost < C^*$
exactly the nodes with $\gcost^*(n) + \hcost_2(n) < C^*$. For such a node,
$\gcost^*(n) + \hcost_1(n) \le \gcost^*(n) + \hcost_2(n) < C^*$, so
\cref{thm:ch04-surely} applied to $\hcost_1$ says that it is expanded there
too. Nodes with $\fcost = C^*$ escape the argument: whether they are expanded
depends on tie-breaking, and $\hcost_2$ may expand a few of them that
$\hcost_1$ skips. Pearl~\cite{pearl1984heuristics} and Dechter and
Pearl~\cite{dechter1985generalized} extend the statement to admissible
heuristics. Their optimality result is a statement about a specific class:
among the admissible algorithms that search the same graph and are
\emph{equally informed}, that is, given the same heuristic and no other
information about the instance, no algorithm expands a strict subset of the
nodes that \astar with a consistent heuristic surely expands. It does not say
that \astar is the fastest possible planner for grids, which it is not
(\cref{sec:ch04-relatives}).
\end{proof}

On grids the theorem orders the heuristics of \cref{tab:ch04-heuristics}: on
an 8-connected grid the octile distance dominates the Euclidean distance,
which dominates the Chebyshev distance, and the self-test of
\code{ch04\_astar.py} checks the theorem on random grids. The theorem also
shows why the maximum of several admissible heuristics is admissible and never
worse than any of them, and it gives the rule of thumb for designing
heuristics: relax the problem (drop the obstacles, drop a constraint), solve
the relaxed problem exactly, and use its cost. That is precisely how the
Manhattan and octile distances arise.

% ---------------------------------------------------------------------
\section{Variants and extensions}
\label{sec:ch04-variants}

\subsection{Tie-breaking}
\label{sec:ch04-tiebreak}

\Cref{thm:ch04-surely} leaves the nodes with $\fcost = C^*$ to the tie-breaking
rule, and near the goal there are many of them: every node on an optimal path
has $\gcost^*(n) + \hcost^*(n) = C^*$, and with an exact heuristic all of them
tie. On an obstacle-free grid the Manhattan distance is exact, and every cell
of the $20 \times 20$ grid in \cref{fig:ch04-tiebreak} lies on some optimal
path from the bottom-left to the top-right corner: all $400$ cells have
$\fcost = 38$. Which of them are expanded is entirely up to the tie-breaking
rule. Preferring the larger $\gcost$ (equivalently, the smaller $\hcost$)
means: among nodes that look equally good, take the one that is furthest along.
The search then behaves like a depth-first dive along one optimal path and
expands $2 \cdot 20 - 1 = 39$ cells. Preferring the smaller $\gcost$, or
breaking ties first-in-first-out without looking at $\gcost$, expands all $400$
cells before the goal is popped. This is why \cref{alg:ch04-astar} uses the key
$(\fcost, -\gcost)$ and why the implementation in
\cref{sec:ch04-implementation} adds a counter as a third component: the counter
keeps the order deterministic and stops \python from comparing nodes when the
first two components tie. In the worked example the rule saves three of 24
expansions; in the last stretch of any search with a good heuristic it saves
far more.\index{tie-breaking!larger g first@larger $\gcost$ first}

\begin{figure}[tb]
  \centering
  \inputfigure{ch04/tiebreak}
  \caption{Tie-breaking on an empty $20 \times 20$ grid with the Manhattan
  heuristic. Every cell has $\fcost = 38 = C^*$. (a) Breaking ties towards the
  larger $\gcost$, \astar expands only the 39 cells of one optimal path.
  (b) Breaking ties towards the smaller $\gcost$, it expands all 400 cells.
  Both runs return a path of cost~38.}
  \label{fig:ch04-tiebreak}
\end{figure}

\subsection{Weighted \astar}
\label{sec:ch04-weighted}

Pohl~\cite{pohl1970heuristic} proposed to inflate the heuristic:
\textbf{weighted \astar}\index{weighted A*}\index{A*!weighted} runs
\cref{alg:ch04-astar} with the key
\begin{equation}
  \fcost_w(n) = \gcost(n) + w\,\hcost(n), \qquad w \ge 1.
  \label{eq:ch04-weighted}
\end{equation}
For $w = 1$ this is \astar; as $w \to \infty$ the $\gcost$ term becomes
irrelevant and the search turns into \emph{greedy best-first search}, which
always expands the node that seems closest to the goal. In between, a larger
$w$ makes the search more greedy: it trusts the heuristic more, expands fewer
nodes and accepts a longer path. The price is bounded.

\begin{theorem}[Bounded suboptimality of weighted \astar]
\label{thm:ch04-weighted}
If $\hcost$ is admissible and a goal is reachable, weighted \astar with
re-opening returns a path of cost at most $w \cdot C^*$.
\end{theorem}
\begin{proof}
\Cref{thm:ch04-invariant} did not use the heuristic at all, so it holds for
the key \cref{eq:ch04-weighted}: when the goal $\gamma'$ is popped, \Open
contains a node $n_i$ of an optimal path $P$ with $\gcost(n_i) =
\gcost^*(n_i)$. As in the proof of \cref{thm:ch04-optimal},
$\hcost(n_i) \le \hcost^*(n_i) \le C^* - \gcost^*(n_i)$, hence
\[
  \fcost_w(n_i) = \gcost^*(n_i) + w\,\hcost(n_i)
    \le \gcost^*(n_i) + w\,\bigl(C^* - \gcost^*(n_i)\bigr)
    \le w\,C^*,
\]
where the last step uses $w \ge 1$ and $\gcost^*(n_i) \ge 0$. The goal was
popped before $n_i$, so $\gcost(\gamma') = \fcost_w(\gamma') \le \fcost_w(n_i)
\le w\,C^*$.
\end{proof}

Note that $w\,\hcost$ is in general not admissible and not consistent, even
when $\hcost$ is, so weighted \astar may re-open nodes and \cref{thm:ch04-consistent}
does not apply. Likhachev, Gordon and Thrun~\cite{likhachev2003ara} showed
that, \emph{when $\hcost$ is consistent}, the bound $w\,C^*$ survives even
when re-opening is switched off and each node is expanded at most once per
iteration; the assumption matters, because with a merely admissible $\hcost$
the argument breaks. That observation is the basis of \arastar in \cref{ch:ch06}, which decreases
$w$ step by step while reusing the search. The focal search of \ecbs in
\cref{ch:ch10} obtains the same guarantee by a different mechanism.

\subsection{An experiment on random grids}
\label{sec:ch04-experiment}

\Cref{fig:ch04-expansions} compares the effort of the searches in this chapter on
random grids with 25\,\% obstacles (script \code{gen\_ch04\_expansions.py},
ten instances per size, start and goal in opposite corners). Three
observations. First, Dijkstra's algorithm expands nearly every free cell
($7\,452$ of $7\,513$ at $n = 100$), while \astar with the Manhattan
heuristic on the 4-connected grid expands $756$ on average: the heuristic
removes about $90\,\%$ of the work. Second, on the 8-connected grid \astar
with the octile heuristic still expands $3\,324$ cells. Diagonal moves make
detours around obstacles cheap, so many cells have $\gcost^* + \hcost <
C^*$; the error $\hcost^* - \hcost$, not the connectivity, is what the search
pays for. Third, weighted \astar with $w = 2$ needs $372$ expansions on the
4-connected and $271$ on the 8-connected grid, while its paths are on average
only $12\,\%$ and $6\,\%$ longer than optimal, far inside the guaranteed factor
of~2; with $w = 1.5$ the average excess is $7\,\%$ and $5\,\%$. Bounded
suboptimality is cheap on these maps, which is why \ecbs (\cref{ch:ch10}) is
so much faster than \cbs.

\begin{figure}[tb]
  \centering
  \inputfigure{ch04/expansions}
  \caption{Left: mean number of expansions (log scale) on random $n \times n$
  grids; Dijkstra's algorithm grows with the number of free cells, \astar grows
  much more slowly, and the weight $w$ cuts the effort again. Right: the mean
  ratio between the cost of the path returned by weighted \astar and the
  optimal cost; the guaranteed bounds are $1.5$ and $2$, the observed ratios
  stay below $1.13$.}
  \label{fig:ch04-expansions}
\end{figure}

\subsection{Other relatives}
\label{sec:ch04-relatives}

\astar has a large family. \lpastar and \dstarlite (\cref{ch:ch05}) reuse the
search when edge costs change; \arastar (\cref{ch:ch06}) is weighted \astar
with a decreasing weight; focal search (\cref{ch:ch10}) keeps a second queue
of nodes whose $\fcost$ is within a factor $w$ of the best and picks among
them by a different criterion. Jump point search~\cite{harabor2011jps}
exploits the symmetry of uniform grids to skip most cells, and safe interval
path planning~\cite{phillips2011sipp} compresses the time dimension of the
search described next. All of them keep the structure of
\cref{alg:ch04-astar}; what changes is the key, the successor function or
what is remembered between searches.

% ---------------------------------------------------------------------
\section{Space-time \astar}
\label{sec:ch04-spacetime}

Every search so far assumed a static world: an obstacle blocks a cell, and it
blocks it for ever. A drone that shares its airspace with other drones faces a
different problem. The cell in front of it may be free now and occupied in three
seconds, when a neighbour flies through it; a corridor may be usable in one
direction only, because a drone is coming the other way. The planner that
resolves such interactions is a \textbf{multi-agent path finding}
(MAPF)\index{MAPF (multi-agent path finding)} algorithm, and the standard
formulation is the one of Stern et al.~\cite{stern2019mapf}, stated in full in
\cref{ch:ch07}: agents move on a shared graph in discrete time steps, each agent
may move to a neighbour or wait, and no two agents may occupy the same vertex at
the same time step or swap places along an edge.
\Cref{ch:ch08,ch:ch09,ch:ch10} build three planners on top of it. All three share one component, and it is an \astar search:
plan for \emph{one} agent, in time, subject to a list of forbidden
$(\text{cell}, \text{time})$ and $(\text{edge}, \text{time})$ pairs. That search
is \textbf{space-time \astar}\index{A*!space-time}\index{space-time A*}, the
subject of this section.

\subsection{The state is a cell and a time}

\begin{definition}[Time-expanded graph and space-time state]
\label{def:ch04-spacetime}
Recall the space-time state and the time-expanded graph of
\cref{def:ch02-space-time-state,def:ch02-time-expanded-graph}; here every
action, move or wait, costs one time step, so time and cost coincide. Let
$G = (V, E)$ be the grid graph of \cref{def:ch02-grid}. A \textbf{space-time
state}\index{space-time state} is a pair $(v, t)$ with $v \in V$ and
$t \in \{0, 1, 2, \dots\}$: the agent is in cell $v$ at time step $t$. From
$(v, t)$ the agent may
either \textbf{wait}\index{wait action}, reaching $(v, t+1)$, or move along an
edge $(v, v') \in E$, reaching $(v', t+1)$. Every action costs one time step. The graph on these states, with these edges,
is the \textbf{time-expanded graph}\index{time-expanded graph} of $G$.
\end{definition}

Three consequences of \cref{def:ch04-spacetime} shape the search.

\emph{The graph is a directed acyclic graph.} Every edge increases $t$ by
exactly one, so no cycle exists, and the cost of a path from $(s, 0)$ to
$(v, t)$ is exactly $t$: $\gcost(v, t) = t$ for every state, no matter which
route reached it. A state therefore has a single possible $\gcost$, and the
first time the search generates it, it is already the cheapest. Relaxation
(line~\ref{alg:ch04-astar:relax} of \cref{alg:ch04-astar}) is replaced by plain
duplicate detection: a state that has been generated is never generated again,
and no node is ever re-opened. The two theorems of
\cref{sec:ch04-properties} apply as they stand, with the trivial consistency
check below.

\emph{Time is the cost.} Cost and time are the same quantity here, which is why
a diagonal move on an 8-connected grid also costs one step rather than
$\sqrt{2}$: a drone that flies diagonally takes one time slot, like every other
manoeuvre. If you need geometric length as well, carry it as a second number and
optimise time first; \cref{ch:ch09} discusses the trade-off, and safe interval
path planning~\cite{phillips2011sipp} avoids the discretisation altogether.

\emph{The state space is infinite unless you bound it.} Waiting is always
allowed, so from every state there is an infinite chain of wait actions. The
time horizon of \cref{sec:ch04-horizon} cuts the chain at a point where nothing
is lost.

\subsection{Constraints and reservation tables}

\begin{definition}[Vertex and edge constraint]
\label{def:ch04-constraints}
In the notation of Stern et al.~\cite{stern2019mapf}, a \textbf{vertex
constraint}\index{constraint!vertex} $\langle a, v, t \rangle$ forbids agent $a$
to occupy cell $v$ at time step $t$; an \textbf{edge
constraint}\index{constraint!edge} $\langle a, u, v, t \rangle$ forbids agent
$a$ to move from $u$ to $v$ between time steps $t$ and $t+1$. A path
$\pi = (\pi_0, \pi_1, \dots, \pi_T)$ of agent $a$, with $\pi_0$ its start and
$\pi_T$ its goal, \textbf{respects} a set of constraints if $\langle a, \pi_t,
t\rangle$ is in none of them for $0 \le t \le T$ and $\langle a, \pi_t,
\pi_{t+1}, t\rangle$ is in none of them for $0 \le t < T$. By the stay-at-goal
convention of \cref{def:ch02-time-indexed-path} the agent still occupies
$\gamma = \pi_T$ at every $t > T$, so a path that respects the constraints must
also avoid $\langle a, \gamma, t\rangle$ for every $t > T$.
\end{definition}

A single-agent search only ever sees the constraints of its own agent, so the
agent index is stripped off before the search starts
(\code{constraints\_for\_agent()} in the chapter code) and the search works with
pairs $(v, t)$ and triples $(u, v, t)$.

Where do the constraints come from? In \cref{ch:ch09} the high level of \cbs
invents them, one per node of its constraint tree, to resolve a collision it has
detected. In prioritized planning (\cref{ch:ch08}) they come from the paths of
the agents that were planned earlier, collected in a \textbf{reservation
table}\index{reservation table}, the data structure Silver introduced for
cooperative pathfinding~\cite{silver2005cooperative}. Reserving a planned path
$\pi$ means three things:
\begin{itemize}
  \item a vertex constraint $(\pi_t, t)$ for every $t$, so that no one flies
        through an occupied cell;
  \item an edge constraint $(\pi_{t+1}, \pi_t, t)$ for every move, that is, the
        \emph{reverse} of the move, so that no one swaps places with $\pi$
        head-on;
  \item a permanent block of the goal cell $\pi_T$ for every $t \ge T$: the
        agent parks there and does not evaporate.
\end{itemize}
The class \code{ReservationTable} in \code{ch04\_astar.py} stores exactly these
three sets and answers the two questions the search asks: is $v$ blocked at time
$t$, and is the move $u \to v$ blocked at time $t$.

\subsection{The goal test: the agent has to stay}
\label{sec:ch04-goalstay}

Reaching the goal is not enough. The plan of an agent ends when it lands, and
the drone then occupies its landing pad for the rest of the episode. If some
other agent is scheduled to fly over that pad at time step~9, arriving there at
step~7 and stopping is not a solution.

\begin{definition}[Goal-stay time]
\label{def:ch04-goalstay}
Let $t_\gamma$ be the largest time step at which the goal $\gamma$ carries a
vertex constraint ($t_\gamma = -1$ if it carries none, $t_\gamma = \infty$ if
another agent is parked on it). A state $(\gamma, t)$ is a
\textbf{goal state}\index{goal!stay condition} if and only if $t > t_\gamma$.
\end{definition}

With this test the returned path $\pi$ can be extended by waiting at $\gamma$ for
ever without violating a constraint, which is what \cref{def:ch04-constraints}
demands of a MAPF solution by the stay-at-goal convention of
\cref{def:ch02-time-indexed-path}. The value $t_\gamma$ is computed once, before the
search, by \code{last\_blocked\_time()}. If $t_\gamma = \infty$ the instance is
unsolvable and the search reports failure without expanding anything.

\begin{pitfall}[Forgetting that the goal must stay free after arrival]
Using the plain test ``$v = \gamma$'' turns a correct plan into a collision:
the drone lands and the next agent flies into it. Using a test that is too
strong is just as bad. A frequent variant is ``$v = \gamma$ and $t \ge H$'',
with $H$ the last constrained time step of the \emph{whole} table; on a map
where some far-away agent is still moving at step~50, every plan is then padded
to 50 steps, and a sub-problem that has a 3-step solution is reported as
needing 50 or, with a tight horizon, as unsolvable. Test the constraints of
\emph{this} goal cell only, as in \cref{def:ch04-goalstay}. In the chapter code
the vertex constraint $\langle a, (2,1), 5\rangle$ on the goal of the mini
example makes the search hop off the goal and come back: the returned path is
$(0,1), (1,1), (2,1), (2,1), (2,1), (2,2), (2,1)$ of cost~6, and it is optimal.
\end{pitfall}

\subsection{Time horizon, completeness and termination}
\label{sec:ch04-horizon}

Because waiting is always possible, the time-expanded graph is infinite, and an
unsolvable instance would make the search run for ever. Cutting the search at an
arbitrary horizon is not safe either: it can report failure on an instance that
has a solution one step later. The following bound is both safe and easy to
compute. Write $H$ for the last time step at which the table blocks anything:
the largest $t$ of a vertex constraint, the largest $t+1$ of an edge constraint,
and the time from which each parked agent occupies its cell ($H = -1$ if the
table is empty). This is what \code{ReservationTable.horizon()} computes. Write
$D$ for the largest number of moves from any cell to $\gamma$ --- among the
cells from which $\gamma$ is reachable at all --- in the grid in which the
parked cells of the table are treated as walls.

\begin{proposition}[A sufficient time horizon]
\label{thm:ch04-horizon}
If a path from $s$ that respects the constraints and satisfies
\cref{def:ch04-goalstay} exists, then one exists that reaches $\gamma$ no later
than time
\begin{equation}
  T_{\max} = H + 1 + D.
  \label{eq:ch04-horizon}
\end{equation}
Space-time \astar restricted to states with $t \le T_{\max}$ is therefore
complete, and it terminates on every instance.
\end{proposition}
\begin{proof}
Let $\pi$ be a solution and suppose it arrives after $T_{\max}$. No vertex or
edge constraint is active at a time step larger than $H$, so from time $H+1$ on
the only forbidden cells are the parked ones; every parked agent has arrived by
time $H$, so from $H+1$ on those cells are blocked for ever. Let
$v = \pi_{H+1}$ be the cell of the agent at time $H+1$. The suffix of $\pi$ from
$H+1$ is a walk from $v$ to $\gamma$ that avoids the parked cells, so $\gamma$ is
reachable from $v$ in the reduced grid and a shortest such route has at most $D$
moves. Following it from time $H+1$ gives a path that respects every constraint,
arrives no later than $H + 1 + D$, and satisfies \cref{def:ch04-goalstay}
because $H \ge t_\gamma$. Its prefix up to $H+1$ is the prefix of $\pi$, which
respects the constraints by assumption. For termination, the set of states with
$t \le T_{\max}$ is finite and each of them is generated at most once, so the
loop ends after finitely many iterations; if it ends without reaching a goal
state, no solution exists.
\end{proof}

\Cref{eq:ch04-horizon} is what \code{default\_horizon()} returns. On the mini
example below, $H = 1$ and $D = 3$, so $T_{\max} = 5$, while the solution needs
3 steps; the bound is loose on purpose, since a loose bound costs nothing when
the search finds a solution early and only matters when it has to prove that
none exists. Supplying a smaller horizon by hand is allowed and useful in a
real-time planner, but it makes the search incomplete: with
\code{max\_time = 2} the mini example is reported unsolvable, with
\code{max\_time = 3} it is solved.

\subsection{The heuristic: static distances, computed once}
\label{sec:ch04-static}

Which heuristic should the search use? The Manhattan distance is admissible on a
4-connected grid, but it ignores the obstacles, and in a constrained re-run of
the same agent that error is paid over and over. The remark at the end of
\cref{sec:ch04-grid-heuristics} said that the exact heuristic $\hcost^*$ is too
expensive for a single query because computing it costs a full search. In
space-time the arithmetic changes completely:

\begin{itemize}
  \item the goal of an agent does not change between the calls, so one backward
        search per agent serves all of them --- in a \cbs run with thousands of
        constraint-tree nodes, that is one search instead of thousands;
  \item the same table serves every time layer, because the distance from $v$ to
        $\gamma$ does not depend on $t$;
  \item on a grid the backward search is a breadth-first search, not a
        Dijkstra search, because every action costs one time step.
\end{itemize}

So let $\hcost(v)$ be the number of moves of a shortest \emph{unconstrained}
path from $v$ to $\gamma$, computed once by a backward breadth-first search
(\code{static\_distances()}), and let $\hcost(v, t) = \hcost(v)$.

\begin{proposition}[The static heuristic is consistent]
\label{thm:ch04-static}
$\hcost$ is admissible and consistent on the time-expanded graph, for any set of
constraints.
\end{proposition}
\begin{proof}
Consistency: for the wait action, $\hcost(v) \le 1 + \hcost(v)$; for a move
$(v, v')$, $\hcost(v) \le 1 + \hcost(v')$, because appending the move to a
shortest route from $v'$ gives a route from $v$. At a goal state $\hcost(\gamma)
= 0$. Consistency implies admissibility by
\cref{thm:ch04-consistent-admissible}; alternatively, note that removing states
and edges from a graph cannot shorten a path, so the unconstrained distance is a
lower bound on the constrained one. Note that $\hcost$ is a lower bound but not
exact: it neither knows the constraints nor the waits they force, which is
exactly the room in which the search does its work.
\end{proof}

A cell with $\hcost(v) = \infty$ cannot reach the goal at all and is never
generated. By \cref{thm:ch04-static} the search inherits
\cref{thm:ch04-consistent}: expansions come out in non-decreasing order of
$\fcost = t + \hcost(v)$, and no state is expanded twice.

\subsection{The algorithm}

\Cref{alg:ch04-spacetime} puts the pieces together. It differs from
\cref{alg:ch04-astar} in four places only: the successor function adds the wait
action and filters on the constraints, the goal test is the one of
\cref{def:ch04-goalstay}, the horizon truncates the search, and relaxation
becomes duplicate detection.

\begin{algorithm}[htb]
\caption{Space-time \astar for one agent under vertex and edge constraints}
\label{alg:ch04-spacetime}
\KwIn{grid graph $G$, start $s$, goal $\gamma$, constraint set (or reservation table) $R$ of this agent, horizon $T_{\max}$}
\KwOut{a time-indexed path $\pi_0, \dots, \pi_T$ of minimum arrival time $T$ that respects $R$, or \emph{failure}}
\Fn{\AstarSpaceTime{$G$, $s$, $\gamma$, $R$, $T_{\max}$}}{
  $\hcost \gets$ static distances to $\gamma$ by a backward breadth-first search in $G$\label{alg:ch04-spacetime:h}\;
  $t_\gamma \gets$ largest $t$ with $(\gamma, t) \in R$ ($-1$ if none, $\infty$ if an agent is parked on $\gamma$)\;
  \lIf{$\hcost(s) = \infty$ or $t_\gamma = \infty$}{\KwRet \emph{failure}\label{alg:ch04-spacetime:hopeless}}
  \Open $\gets$ priority queue holding $(s, 0)$ with key $(\hcost(s),\, 0)$; $\mathit{Gen} \gets \{(s, 0)\}$; parent$(s, 0) \gets$ \KwNil\;
  \While{\Open $\ne \emptyset$}{
    $(v, t) \gets$ pop the entry with the smallest key $(\fcost, -t)$ from \Open\label{alg:ch04-spacetime:pop}\;
    \If{$v = \gamma$ \KwAnd $t > t_\gamma$}{\KwRet \AstarPath{parent, $(v,t)$}\label{alg:ch04-spacetime:goal}\;}
    \If{$t = T_{\max}$}{\KwContinue\tcp*{horizon reached: no successors}\label{alg:ch04-spacetime:horizon}}
    \ForEach{$v' \in \{v\} \cup \Succ(v)$}{\tcp*[h]{wait first, then the moves}\label{alg:ch04-spacetime:succ}\;
      \lIf{$\hcost(v') = \infty$}{\KwContinue\label{alg:ch04-spacetime:dead}}
      \lIf{$(v', t+1) \in R$ \KwOr $(v, v', t) \in R$}{\KwContinue\label{alg:ch04-spacetime:filter}}
      \lIf{$(v', t+1) \in \mathit{Gen}$}{\KwContinue\tcp*{$\gcost = t+1$ always: the first is the best}\label{alg:ch04-spacetime:dup}}
      $\mathit{Gen} \gets \mathit{Gen} \cup \{(v', t+1)\}$; parent$(v', t+1) \gets (v, t)$\;
      push $(v', t+1)$ into \Open with key $(t + 1 + \hcost(v'),\, -(t+1))$\label{alg:ch04-spacetime:push}\;
    }
  }
  \KwRet \emph{failure}\;
}
\end{algorithm}

Line~\ref{alg:ch04-spacetime:h} is the backward breadth-first search of
\cref{sec:ch04-static}; in a planner that calls the search many times for the
same agent it is hoisted out and computed once.
Line~\ref{alg:ch04-spacetime:hopeless} catches the two hopeless cases before any
work is done. Line~\ref{alg:ch04-spacetime:succ} generates the wait action
first, so that among states with equal keys the wait is popped first; this is a
choice of tie-breaking, not of correctness.
Line~\ref{alg:ch04-spacetime:dead} drops cells from which the goal cannot be
reached at all, and line~\ref{alg:ch04-spacetime:filter} is the whole of the
constraint handling: one lookup in a hash set per candidate.
Line~\ref{alg:ch04-spacetime:dup} is duplicate detection; because
$\gcost(v', t+1) = t+1$ for every route, a second parent could never be better.
The set $\mathit{Gen}$ of generated states plays the role of both \Open and \Closed of
\cref{alg:ch04-astar}. Note that the search minimises the \emph{arrival time}
$T$; if an agent must also stand still at its goal until some later step, that
is a property of the plan, not of the search.

\subsection{A worked mini-example}

\begin{example}[One vertex constraint, one forced wait]
\label{ex:ch04-spacetime}
Take the empty $3 \times 3$ grid, 4-connected, with $s = (0, 1)$ and
$\gamma = (2, 1)$, and the single vertex constraint $\langle a, (1,1), 1\rangle$:
another drone crosses the middle cell at time step~1. The instance is
\code{spacetime\_example()} in \code{ch04\_astar.py}. Without the constraint the
search expands the three states $((0,1),0)$, $((1,1),1)$, $((2,1),2)$ and
returns the path $(0,1), (1,1), (2,1)$ of cost~2: with an exact heuristic and
$\fcost = 2$ everywhere, the tie-breaking rule walks straight down the optimal
path. With the constraint, the state $((1,1),1)$ no longer exists, the search
expands four states and returns
\[
  (0,1),\ (0,1),\ (1,1),\ (2,1),
\]
a path of cost~3 with one wait at the start. \Cref{fig:ch04-spacetime} shows the
two paths in the time-expanded graph and \cref{tab:ch04-spacetime} the trace.
\end{example}

\begin{figure}[tb]
  \centering
  \inputfigure{ch04/spacetime}
  \caption{\Cref{ex:ch04-spacetime}. Left: the $3\times3$ grid; the hatched cell
  $(1,1)$ is forbidden at $t = 1$. Right: the states $(v, t)$ of the middle row
  for $t = 0, \dots, 3$. Horizontal edges are wait actions, diagonal edges are
  moves, and every edge costs one time step. The crossed-out state $((1,1),1)$
  is the one removed by the vertex constraint; it lies on the dashed path of
  cost~2, which is therefore no longer available. Grey states are the four the
  search expands, with their $\fcost = t + \hcost(v)$; the cheapest remaining
  path (solid) waits once and costs~3. The states of the top and bottom rows,
  such as $((0,0),1)$ and $((0,2),1)$ in \cref{tab:ch04-spacetime}, exist too
  but are never expanded, so the drawing omits them.}
  \label{fig:ch04-spacetime}
\end{figure}

\begin{table}[tb]
  \centering\footnotesize
  \caption{Trace of \cref{alg:ch04-spacetime} on \cref{ex:ch04-spacetime}. Note
  that $\gcost$ is always the time $t$, that $\fcost$ never decreases
  (\cref{thm:ch04-consistent}(ii) applies, since $\hcost$ is consistent), and
  that the state $((1,1),1)$ never enters \Open: the constraint filter of
  line~\ref{alg:ch04-spacetime:filter} removes it at generation time.}
  \label{tab:ch04-spacetime}
  \begin{tabular}{@{}rlrrrp{5.1cm}@{}}
  \toprule
  Step & Expanded $(v,t)$ & $\gcost$ & $\hcost$ & $\fcost$ & \Open after the step (state$:\fcost$)\\
  \midrule
  1 & $((0,1),0)$ & 0 & 2 & 2 & $((0,0),1){:}4$, $((0,1),1){:}3$, $((0,2),1){:}4$\\
  2 & $((0,1),1)$ & 1 & 2 & 3 & $((1,1),2){:}3$, $((0,1),2){:}4$, $((0,0),1){:}4$, $((0,2),1){:}4$, $((0,0),2){:}5$, $((0,2),2){:}5$\\
  3 & $((1,1),2)$ & 2 & 1 & 3 & $((2,1),3){:}3$, $((1,1),3){:}4$, $((0,1),3){:}5$, $((1,0),3){:}5$, $((1,2),3){:}5$, and the five entries of step~2 that are still open\\
  4 & $((2,1),3)$ & 3 & 0 & 3 & goal state, $t = 3 > t_\gamma = -1$: return the path\\
  \bottomrule
  \end{tabular}
\end{table}

Two variations, both in the self-test, show the other two mechanisms at work.
Adding the edge constraint $\langle a, (0,1), (1,1), 1\rangle$ on top of the
vertex constraint forbids the move that the path above makes at $t = 1$; the
search then waits twice and returns $(0,1), (0,1), (0,1), (1,1), (2,1)$ of
cost~4 after five expansions. Replacing the constraints by the single vertex
constraint $\langle a, (2,1), 5\rangle$ on the goal makes $t_\gamma = 5$: the
search reaches $(2,1)$ at $t = 2$ but may not stop there, so it waits, steps
aside to $(2,2)$ at $t = 5$ and returns at $t = 6$, for a path of cost~6.

\begin{example}[Two drones in a corridor with a passing bay]
\label{ex:ch04-corridor}
A $5 \times 2$ map has a corridor along $y = 0$ and a single free cell $(2, 1)$
above it, a passing bay. Agent~1 flies from $(0,0)$ to $(4,0)$ and is planned
first; its path $(0,0), (1,0), (2,0), (3,0), (4,0)$ goes into the reservation
table, together with the reversed edges and the permanent block of $(4,0)$ from
$t = 4$. Agent~2 has to fly from $(3,0)$ to $(0,0)$, which on the empty corridor takes
3 steps. Under the table, space-time \astar returns
\[
  (3,0),\ (2,0),\ (2,1),\ (2,0),\ (1,0),\ (0,0),
\]
of cost~5: agent~2 hides in the bay at $t = 2$, lets agent~1 pass, and drops
back into the corridor behind it. No sequence of waits in the corridor itself
would work, because the reversed-edge constraints forbid the head-on swap. If
agent~2 instead starts at $(4,0)$, no path exists at all --- agent~1 parks on
that cell --- and the search proves it in finite time, six expansions, thanks to
\cref{thm:ch04-horizon}: here $H = 4$ and $D = 3$, so $T_{\max} = 8$. Every
number in this example is asserted by the self-test of
\code{ch04\_astar.py}. That failure is not a bug in the search but the known
incompleteness of prioritized planning, and it is the reason \cref{ch:ch09}
searches over constraints instead of fixing an order.
\end{example}

% ---------------------------------------------------------------------
\section{Implementation notes}
\label{sec:ch04-implementation}

\Cref{lst:ch04-astar} is the main loop of \code{astar()} in
\code{code/ch04\_astar.py}, the function that produced every number in this
chapter. It is \cref{alg:ch04-astar} line by line, and the rest of this section
explains the six decisions hidden in it.

\begin{lstlisting}[caption={The main loop of \astar (from \code{code/ch04\_astar.py}): a lazy binary heap, a dictionary of $\gcost$ values, a closed set and parent pointers. \code{key()} builds the priority tuple $(\fcost, -\gcost, \mathit{counter})$; \code{EPS} is $10^{-9}$.},label={lst:ch04-astar}]
    g = {start: 0.0}
    parent = {start: None}
    closed = set()
    expansions = []
    reopenings = 0
    open_heap = [(key(weight * heuristic(start), 0.0), 0.0, start)]
    while open_heap:
        _, g_entry, n = heapq.heappop(open_heap)
        if g_entry > g[n]:               # stale: a cheaper path was found
            continue
        h_n = heuristic(n)
        expansions.append(Expansion(n, g[n], h_n, g[n] + weight * h_n))
        if is_goal(n):
            snapshot = live_open()
            if record_open:
                expansions[-1].open_after = snapshot
            return SearchResult(_reconstruct(parent, n), g[n], expansions,
                                closed, {m for (m, _) in snapshot}, g,
                                reopenings)
        closed.add(n)
        for n2, c in successors(n):
            g2 = g[n] + c
            if g2 < g.get(n2, INF) - EPS:  # relaxation (EPS: fp noise)
                if n2 in closed:
                    if not reopen:
                        continue
                    closed.discard(n2)   # re-open (never for consistent h)
                    reopenings += 1
                g[n2] = g2
                parent[n2] = n
                f2 = g2 + weight * heuristic(n2)
                heapq.heappush(open_heap, (key(f2, g2), g2, n2))
        if record_open:
            expansions[-1].open_after = live_open()
    return SearchResult(None, INF, expansions, closed, set(), g, reopenings)
\end{lstlisting}

\paragraph{Hashing the state.}
The state is the dictionary key, so it must be hashable and cheap to hash. A
grid cell is the tuple \code{(x, y)} of two \code{int}s; a space-time state is
the nested tuple \code{((x, y), t)}. Both hash in constant time and compare by
value, which is what the sets \code{closed} and \code{parent} need. Do not use a
NumPy array or a list as a state: they are unhashable, and a class without
\code{\_\_hash\_\_} and \code{\_\_eq\_\_} would compare by identity, so two
states describing the same cell would look different and the search would
degenerate into a tree search.

\paragraph{The priority tuple and why it needs a counter.}
The heap entries are \code{(key, g, node)} with
\code{key = (f, -g, counter)}. \python compares tuples lexicographically, so the
first component orders by $\fcost$, the second breaks ties towards the larger
$\gcost$ (\cref{sec:ch04-tiebreak}), and the third, a value from
\code{itertools.count()}, breaks the remaining ties first-in-first-out. The
counter is not cosmetic. Without it, two entries with identical $\fcost$ and
$\gcost$ would make \code{heapq} compare the next component, the node itself;
for grid cells that silently sorts by coordinate, which is a bias nobody
intended, and for a state that is not comparable at all, such as a
\code{dataclass} without an ordering, it raises \code{TypeError} in the middle
of a long run. The counter guarantees that no two keys are ever equal, so the
node is never reached by the comparison.

\paragraph{Parents and path reconstruction.}
A dictionary \code{parent[node] = predecessor} costs one pointer per generated
node and turns path extraction into a walk from the goal back to the start
(\code{\_reconstruct()}), which is then reversed. Storing whole paths in the
queue instead is the classic beginner's implementation; it multiplies the memory
by the path length and makes every push copy a list.

\paragraph{Lazy deletion instead of decrease-key.}
When the $\gcost$ of a node improves, the tidy operation is
\emph{decrease-key}, which \code{heapq} does not offer. The implementation
pushes a second entry instead and lets the stale one age in the heap; the test
\code{g\_entry > g[n]} on the pop discards it. The heap therefore holds at most
one entry per relaxation, so its size is bounded by the number of edges, and the
cost of a pop stays $\bigO{\log \abs{E}} = \bigO{\log \abs{V}}$ on a grid. The
space-time search does not need the mechanism at all, because $\gcost$ is the
time and can never improve.

\paragraph{Memory.}
Memory, not time, is what stops \astar on large problems. Every generated node
occupies an entry in \code{g}, in \code{parent} and, until it is popped, in the
heap. On a $1000 \times 1000$ grid with 4-connectivity, a search that touches
every cell holds a few million small \python objects, hundreds of megabytes.
This is why \cref{ch:ch06} discusses anytime and memory-bounded variants, and
why the space-time search bounds its state space with \cref{eq:ch04-horizon}
rather than letting the wait actions run free.

\paragraph{Floating point: \code{round(f, 9)} and \code{EPS}.}
On an 8-connected grid the costs are $1$ and $\sqrt{2}$, so $\gcost$ is a sum of
irrational numbers and two routes of equal true cost can differ in the last bit.
Two places react badly to that. The relaxation test
\code{g2 < g.get(n2, INF) - EPS} would otherwise accept an ``improvement'' of
$10^{-16}$, re-open the node, push it again and, with an unlucky cycle of
rounding, loop; subtracting \code{EPS} makes the test insist on a real
improvement. The key uses \code{round(f, 9)} so that two $\fcost$ values that
are equal in exact arithmetic are also equal in the tuple, which lets the
$-\gcost$ component do the tie-breaking it was put there for --- without the
rounding, the tie-breaking experiment of \cref{sec:ch04-tiebreak} is decided by
floating-point noise. The self-test relies on this: it checks that the $\fcost$
sequence is non-decreasing and that no node is re-opened with a consistent
heuristic, and both assertions fail without the guards.

\Cref{lst:ch04-spacetime} shows the corresponding heart of
\code{space\_time\_astar()}: the goal test of \cref{def:ch04-goalstay}, the wait
action, the constraint filter and the duplicate test, in the same order as
lines~\ref{alg:ch04-spacetime:goal}--\ref{alg:ch04-spacetime:push} of
\cref{alg:ch04-spacetime}.

\begin{lstlisting}[caption={The expansion step of space-time \astar (from \code{code/ch04\_astar.py}): the goal-stay test, the wait action, the constraint filter and duplicate detection. \code{h} is the static distance table and \code{table} the reservation table. The \code{closed} set is recorded only for the traces and the figures; correctness needs only \code{parent}.},label={lst:ch04-spacetime}]
    while heap:
        _, (v, t) = heapq.heappop(heap)  # every state is pushed exactly once
        closed.add((v, t))
        expansions.append(Expansion((v, t), t, h[v], t + h[v]))
        if v == goal and t > t_goal:                 # stay-at-goal is safe
            path = [s[0] for s in _reconstruct(parent, (v, t))]
            if record_open:
                expansions[-1].open_after = snapshot()
            return SearchResult(path, float(t), expansions, closed,
                                {s for (_, s) in heap})
        if t < max_time:
            for v2 in [v] + [n for n, _ in succ(v)]:  # wait first, then moves
                if v2 not in h:                       # cannot reach the goal
                    continue
                if (table.vertex_blocked(v2, t + 1)
                        or table.edge_blocked(v, v2, t)):
                    continue
                s2 = (v2, t + 1)
                if s2 in parent:                     # g = t + 1: first is best
                    continue
                parent[s2] = (v, t)
                heapq.heappush(heap, (key(*s2), s2))
\end{lstlisting}

\begin{pitfall}[Testing for the goal when the node is generated]
It is tempting to check ``is this successor the goal?'' inside the successor
loop and to return at once. The path that comes back is then the first one
found, not the cheapest. In \cref{ex:ch04-grid} the goal $(5,2)$ is generated in
step~20 with $\gcost = 13$, which happens to be optimal, but that is luck. On
the four-node graph of \cref{sec:ch04-properties} with the inconsistent
heuristic, the goal is generated by the first expansion of $b$ with
$\gcost = 5$ and only later improved to $4$: a search that returns at
generation time reports a path 25\,\% too expensive. The
proof of \cref{thm:ch04-optimal} breaks at exactly this point --- it compares
$\fcost(\gamma')$ with the keys of the entries \emph{in} \Open when $\gamma'$ is
\emph{popped}, and a generated node has not been popped. Test the goal after the
pop, on line~\ref{alg:ch04-astar:goal}. The one search in this book that may
test at generation time is breadth-first search on unit-cost graphs, where all
paths to a node found at the same depth cost the same.
\end{pitfall}

% ---------------------------------------------------------------------
\section{Where this fits in the drone system}
\label{sec:ch04-drone}

\begin{dronebox}
\astar is the single most reused component of the hybrid planner of
\cref{ch:ch24}. It appears in three roles.

\textbf{1. The low-level search of the multi-agent planner.} \cbs
(\cref{ch:ch09}) and \ecbs (\cref{ch:ch10}) split a collision into constraints
and call, at every node of their constraint tree, a single-agent search for
\emph{one} drone under \emph{that node's} constraints. That call is
\cref{alg:ch04-spacetime}, with the constraints of
\cref{def:ch04-constraints} and the goal test of \cref{def:ch04-goalstay}.
Prioritized planning (\cref{ch:ch08}) makes the same call with a reservation
table instead. A run on thirty drones makes thousands of these calls, so the
static distance table of \cref{sec:ch04-static} is computed once per drone and
reused, and every constant factor in \cref{lst:ch04-astar} is multiplied by a
four-digit number.

\textbf{2. The fallback of the replanning layer.} The reactive layer of
\cref{ch:ch24} --- optimal reciprocal collision avoidance (\orca,
\cref{ch:ch13}) and the dynamic window approach (DWA, \cref{ch:ch14}) --- handles the
seconds-scale avoidance of an intruder locally. When a manoeuvre has pushed a
drone so far off its nominal route that the local controller cannot rejoin it,
the replanning layer asks for a new global path from the current cell, and the
question goes back to \astar --- in space-time again, with the other drones'
remaining plans reserved, so that the repair does not create a new collision.

\textbf{3. The baseline for the incremental planners.} \lpastar and \dstarlite
(\cref{ch:ch05}) repair the previous search instead of restarting it, which is
the right answer when the map changes a little and often; they are \astar plus
bookkeeping, and their correctness arguments are the ones proved here. Use
plain \astar when the change is large or rare, and \cref{ch:ch05} when a drone
discovers obstacles as it flies.

This chapter is the Week~1 milestone of the study plan (\cref{ch:appA}), whose
coding exercise --- grid \astar, time as a state variable, a picture of the open
and closed sets --- is \cref{exr:ch04-coding} below.
\end{dronebox}

% ---------------------------------------------------------------------
\begin{summary}
\begin{itemize}
  \item \astar is best-first search on $\fcost = \gcost + \hcost$: the cost so
        far plus an estimate of the cost to go. With $\hcost \equiv 0$ it is
        Dijkstra's algorithm.
  \item $\hcost$ is \emph{admissible} if $\hcost \le \hcost^*$ everywhere, and
        \emph{consistent} if $\hcost(\gamma) = 0$ and $\hcost(n) \le c(n,n') +
        \hcost(n')$ on every edge. Consistency implies admissibility; the
        converse fails.
  \item With an admissible $\hcost$ and re-opening enabled, the first goal
        popped is optimal (\cref{thm:ch04-optimal}). With a consistent $\hcost$,
        no node is expanded twice and the $\fcost$ of the expansions never
        decreases (\cref{thm:ch04-consistent}); the nodes with $\gcost^* +
        \hcost < C^*$ are exactly the ones that must be expanded
        (\cref{thm:ch04-surely}). Both theorems assume a finite graph,
        non-negative costs and the goal test on the pop.
  \item On grids: Manhattan for 4-connected maps, octile for 8-connected maps,
        Euclidean and Chebyshev when nothing better is available. The Manhattan
        distance is \emph{not} admissible on an 8-connected grid and silently
        returns longer paths (\cref{tab:ch04-heuristics}).
  \item Ties in $\fcost$ decide most of the work near the goal: prefer the
        larger $\gcost$, and add a counter to keep the comparison total.
  \item Weighted \astar with $\fcost_w = \gcost + w\hcost$ returns a path of cost
        at most $w\,C^*$ and is usually far better than the bound.
  \item Space-time \astar searches over states $(v, t)$ with a wait action, one
        time step per action, vertex constraints $\langle a,v,t\rangle$ and edge
        constraints $\langle a,u,v,t\rangle$, a goal test that requires the goal
        to stay free ($t > t_\gamma$), a horizon $T_{\max} = H + 1 + D$ and a
        static-distance heuristic computed once per goal. It is the low-level
        search of every multi-agent planner in this book.
\end{itemize}
\end{summary}

\paragraph*{Further reading.}
\astar was introduced by Hart, Nilsson and Raphael~\cite{hart1968formal}, whose
correction note~\cite{hart1972correction} fixes the definition of the ``more
informed'' relation used in \cref{thm:ch04-dominance}.
Pearl~\cite{pearl1984heuristics} is the reference book on heuristic search and
contains the sharpest form of the optimality results, including the
complexity discussion of \cref{sec:ch04-properties}; Dechter and
Pearl~\cite{dechter1985generalized} settle exactly in which class \astar is
optimal. Russell and Norvig~\cite{russell2020aima} give the textbook treatment
that the proofs above follow, and Pohl~\cite{pohl1970heuristic} the original
weighted search. For grids, jump point search~\cite{harabor2011jps} is the
standard way to exploit the symmetry of uniform maps, and safe interval path
planning~\cite{phillips2011sipp} replaces the time layers of
\cref{sec:ch04-spacetime} by intervals, which pays off when the horizon is long
and the constraints are sparse. Silver~\cite{silver2005cooperative} introduced
the reservation table and cooperative pathfinding in the setting that
\cref{ch:ch08} formalises, and Stern et al.~\cite{stern2019mapf} fix the MAPF
terminology, including the vertex and edge constraints of
\cref{def:ch04-constraints}, used throughout \cref{ch:ch08,ch:ch09,ch:ch10}.

% ---------------------------------------------------------------------
\section{Exercises}
\label{sec:ch04-exercises}

\begin{exercise}[\difficulty{1}]
\label{exr:ch04-trace}
Continue \cref{tab:ch04-trace} by hand for two more steps of \cref{ex:ch04-grid}.
Give, for steps~11 and~12, the expanded cell with its $\gcost$, $\hcost$ and
$\fcost$, and the contents of \Open after the step with the current $\fcost$ of
each cell. Then answer: why is $(0,4)$, which has been in \Open since step~8,
still not expanded?
\end{exercise}

\begin{exercise}[\difficulty{1}]
\label{exr:ch04-fourheuristics}
On the grid of \cref{ex:ch04-grid} with $\gamma = (5,2)$, compute all four
heuristics of \cref{eq:ch04-manhattan}--\cref{eq:ch04-chebyshev} at the cells
$(1,4)$ and $(3,1)$. (a) Order them by dominance and say which is the most
informed. (b) Which of them are admissible for the 4-connected instance, and
which for the 8-connected one? (c) The true costs to go on the
4-connected instance are $\hcost^*(1,4) = 8$ and $\hcost^*(3,1) = 5$. By how
much does the best of the four underestimate, and which obstacle causes the
error?
\end{exercise}

\begin{exercise}[\difficulty{2}]
\label{exr:ch04-octile}
Prove that the octile distance \cref{eq:ch04-octile} is consistent on an
8-connected grid whose moves obey the corner-cutting rule of
\cref{def:ch02-grid}. Show first that $\hcost_{\mathrm{O}}$ is the exact cost on
an obstacle-free grid, then verify \cref{eq:ch04-consistent} on a straight and
on a diagonal move directly from the formula (four cases, depending on whether
the move changes $\max(d_x,d_y)$ or $\min(d_x,d_y)$), and finally explain in one
sentence why forbidding some moves cannot break consistency.
\end{exercise}

\begin{exercise}[\difficulty{2}]
\label{exr:ch04-wbound}
\Cref{thm:ch04-weighted} assumes that closed nodes may be re-opened. Prove the
same bound $w\,C^*$ for the variant that never re-opens a closed node, assuming
that $\hcost$ is consistent. Hint: show first that every expanded node $n$
satisfies $\gcost(n) \le w\,\gcost^*(n)$ by induction along a cheapest path to
$n$, then repeat the argument of \cref{thm:ch04-weighted} with that inequality
in place of \cref{thm:ch04-invariant}. Where exactly does your proof use
consistency, and what goes wrong when $\hcost$ is only admissible?
\end{exercise}

\begin{exercise}[\difficulty{2}]
\label{exr:ch04-max}
Let $\hcost_1$ and $\hcost_2$ be admissible and let $\hcost(n) =
\max(\hcost_1(n), \hcost_2(n))$. (a) Show that $\hcost$ is admissible, and that
it is consistent when both $\hcost_1$ and $\hcost_2$ are. (b) Show that
$\hcost$ dominates both, and state what \cref{thm:ch04-dominance} then predicts
about the number of expansions. (c) Give two admissible heuristics on the
4-connected grid of \cref{ex:ch04-grid} whose maximum is strictly better than
both, and explain why one does not simply always use the maximum of every
heuristic one can think of.
\end{exercise}

\begin{exercise}[\difficulty{2}]
\label{exr:ch04-lattice}
\Cref{ch:ch02} defines 6-, 18- and 26-connected lattices in three dimensions,
with move costs $1$, $\sqrt{2}$ and $\sqrt{3}$ respectively for moves that
change one, two or three coordinates. Give an admissible and consistent
heuristic for each connectivity, exact on an obstacle-free lattice, in the style
of \cref{eq:ch04-octile}. Hint: sort $d_x, d_y, d_z$ so that $d_{(1)} \ge
d_{(2)} \ge d_{(3)}$ and count how many moves of each kind an optimal
obstacle-free route uses. Then say which of them is admissible on which of the
other two lattices, in the manner of \cref{tab:ch04-heuristics}.
\end{exercise}

\begin{exercise}[\difficulty{2}]
\label{exr:ch04-manhattan8}
Construct a grid of at most $4 \times 4$ cells, an 8-connected instance and a
start/goal pair for which \astar with the Manhattan heuristic returns a path
that is strictly longer than the optimum, and give both costs. Then explain,
using \cref{thm:ch04-optimal}, which step of the optimality proof the Manhattan
heuristic invalidates, and check your instance with
\code{astar\_grid(grid, s, z, "manhattan", 8)} against
\code{dijkstra\_grid(grid, s, 8)}.
\end{exercise}

\begin{exercise}[\difficulty{3} Coding]
\label{exr:ch04-coding}
This is the Week~1 coding exercise of the study plan (\cref{ch:appA}); it
continues \cref{exr:ch03-coding}. Implement \astar yourself, in the API of
\code{code/ch04\_astar.py}, and do not read that file until you are stuck.
(a) Write \code{astar(start, goal, successors, heuristic)} with a lazy heap, a
closed set, the priority tuple $(\fcost, -\gcost, \mathit{counter})$ and parent
pointers, and \code{astar\_grid()} on top of it for 4- and 8-connected grids.
Reproduce the numbers of \cref{ex:ch04-grid}: cost~13, 21 expansions, and the
first ten rows of \cref{tab:ch04-trace}. (b) Add \code{record\_open=True} and
draw the open and closed sets of your search, in the style of
\cref{fig:ch04-example}, after every expansion; watch the frontier of the
Manhattan search walk into the dead end of the column $x = 3$. (c) Add time as
a state variable: implement \code{space\_time\_astar(grid, start, goal,
constraints, max\_time)} with the wait action, the constraint filter of
\cref{def:ch04-constraints}, the goal test of \cref{def:ch04-goalstay} and the
horizon of \cref{eq:ch04-horizon}. Check it against \cref{ex:ch04-spacetime} and
\cref{ex:ch04-corridor}. (d) Verify on 50 random $12 \times 12$ grids that your
\astar returns the same cost as your Dijkstra of \cref{ch:ch03}, that no node is
re-opened when the heuristic is consistent, and that the space-time cost without
constraints equals the static distance.
\end{exercise}

\begin{exercise}[\difficulty{3}]
\label{exr:ch04-spacetime}
Take \cref{ex:ch04-spacetime} and add, by hand, the vertex constraint
$\langle a, (1,1), 2\rangle$ to the existing $\langle a, (1,1), 1\rangle$.
(a) Predict the cost of the new optimal path, the number of waits and the route.
(b) Now replace both by the single edge constraint
$\langle a, (1,1), (2,1), 1\rangle$ and predict again; then say what would
change if its time index were $2$ instead of $1$. (c) Check both with
\code{space\_time\_astar()} and compare the expansion counts with the four
expansions of \cref{tab:ch04-spacetime}. (d) Which of the two instances would
become unsolvable if the horizon were set to \code{max\_time = 3}, and why does
that not contradict \cref{thm:ch04-horizon}?
\end{exercise}

\begin{exercise}[\difficulty{3}]
\label{exr:ch04-horizon}
\Cref{thm:ch04-horizon} bounds the arrival time by $T_{\max} = H + 1 + D$.
(a) Show by an example that the term $D$ cannot be dropped, that is, that
$H + 1$ alone is not a sufficient horizon. (b) Show that the bound is not tight
in general by giving an instance whose optimal arrival time is far below
$T_{\max}$. (c) In \cref{ex:ch04-corridor}, compute $H$, $D$ and $T_{\max}$ for
agent~2 and compare with the arrival time~5 that the search returns.
(d) A colleague proposes to drop the parked cells from the computation of $D$,
using the eccentricity of the goal in the raw grid instead. Show that the
resulting horizon can be too small, and describe the failure the planner would
report.
\end{exercise}
